\section{Function-Space Overlap: Neural Operators and Mergeable Sketches}\label{app:operator-overlap}

Section~\ref{sec:fno-primer} introduced neural operators as reusable
function-space maps, and Section~\ref{sec:mergeable-sketches} introduced
mergeable sketches as encode/merge/decode summaries. This appendix explains
the overlap in the least mysterious terms: a mergeable summary builds a state
on each shard, merges states up the tree, and reads the answer from the root
state. C-TreePO uses the same state-level shape when $g$ builds and merges the
state and $f$ reads it out. The point is not that every neural operator is a
mergeable sketch, or that every classical sketch is implemented by a
transformer. The point is that there is a named state-level intersection, and
the local laws certify membership in it.
Appendix~\ref{app:kovachki-finite-dimensionalization} records the
separate Kovachki Lemma~21/22 finite-dimensionalization route; this
appendix uses that route only through the resulting approximation and
local-law subspace interfaces.

\subsection{Notation Map}

Table~\ref{tab:ctreepo-neural-operator-notation} adapts the notation
dictionary role of \citet[Table~10]{Kovachki23NeuralOperator} to the
objects used in this paper. The key shift is that our ``domain'' is not a
PDE mesh alone: it is the index set of the realized call currently being
processed, whether that call is a raw leaf, a pair of child summaries, or
an on-range summary being re-used.

\begin{table}[htbp]
\centering
\footnotesize
\setlength{\tabcolsep}{3pt}
\begin{tabular}{p{0.18\linewidth}p{0.34\linewidth}p{0.40\linewidth}}
\toprule
Neural-operator notation & Meaning in C-TreePO & Examples in this paper \\
\midrule
$D$, $x\in D$ &
Index domain and locations inside one operator call. &
Markov token positions; HLL stream/register coordinates; manifesto character or quasi-sentence offsets. \\
$a\in\mathcal{A}$ &
Input function or state presented to the operator. &
A leaf block, concatenated child summaries, or an on-range summary; Markov observed-token sequence; HLL stream/register state; RILE span plus rubric prompt. \\
$u\in\mathcal{U}$ &
Target output function, summary state, or readout-facing state. &
Markov $(\text{count},\text{first},\text{last})$ state or root count; HLL register vector or estimate; RILE evidence summary or score vector. \\
$D_j$ &
Discretization/refinement at which the same operator family is evaluated. &
Chunk size and leaf geometry; tokenizer/grid length for FNO or transformer calls; HLL precision as the register-state resolution. \\
$\mathcal{G}^{\dagger}:\mathcal{A}\to\mathcal{U}$ &
Ideal target operator before approximation. &
Exact theorem-backed summarizer; hand-designed mergeable sketch; oracle-preserving Markov boundary-state merger. \\
$\mathcal{G}_{\theta}$, $g$ &
Realized learned operator used by the tree. &
1D FNO leaf/merge network; learned HLL-style merge/readout; prompted LLM leaf and merge operator. \\
$\mu$ &
Distribution over inputs or realized calls. &
Synthetic Markov DGP; HLL benchmark stream distribution; Manifesto corpus and sampled audit-node distribution. \\
$v_t$, $d_v$, $P,Q,K,W,\sigma$ &
Hidden representation dimension and equation-(6) layer components. &
FNO latent state and Fourier kernel; transformer attention/kernel layer plus feed-forward maps; MLP merge/readout adapters. \\
$K_{\mathrm{call}}$ &
Compact realized-call set on which approximation is required. &
The finite/compact collection of leaf calls, merge calls, and on-range re-summary calls generated by one audited tree run. \\
\bottomrule
\end{tabular}
\caption{C-TreePO notation map for the neural-operator layer, paralleling the operator-learning notation dictionary of \citet[Table~10]{Kovachki23NeuralOperator} while making the paper's three examples explicit.}
\label{tab:ctreepo-neural-operator-notation}
\end{table}

\subsection{The Diagram}

Fix an ambient operator space
\[
\mathrm{Op}(\mathcal{A},\mathcal{U})
  := \{\,\mathcal{G}:\mathcal{A}\to\mathcal{U}\,\}.
\]
The formalized diagram is:
\[
\begin{tikzcd}[column sep=large,row sep=large]
\text{finite transformer encoders}
  \arrow[r, hook]
  & \mathrm{Eq6NO}
  \arrow[r, hook]
  & \mathrm{NO}
  \arrow[d, "\cap\mathrm{MS}" description] \\
\text{mergeable sketches}
  \arrow[r, hook]
  & \mathrm{ExactLL}(T,\fstar)
  \arrow[r, hook]
  & \mathrm{ApproxLL}(T,\fstar;\varepsilon)
\end{tikzcd}
\]
Here $\mathrm{Eq6NO}$ is the equation-(6) class from
Equation~\eqref{eq:equation6-no}, $\mathrm{MS}$ is the class of
operators induced by mergeable sketches with exact sketch-local
witnesses, $\mathrm{ExactLL}$ is the exact C1/C2/C3 subspace, and
$\mathrm{ApproxLL}$ is its budget-relaxed version. The right vertical
label is the certified overlap:
\[
\mathrm{MNO}_{\mathcal{N},\fstar}
  :=
  \mathcal{N}\cap\mathrm{MS},
\]
where $\mathcal{N}$ is the chosen neural-operator class.
Read this as a set-membership statement, not as an automatic guarantee about
trained models. Neural-operator architecture gives a representation family.
Mergeable-summary structure gives build/merge/readout witnesses. C1/C3/C2
are the checks that a realized operator is actually using the state in the
required way.

\subsection{Agarwal Notation and the \texorpdfstring{$\varepsilon$}{epsilon}-Fixed Interface}

For the comparison to \citet{Agarwal2013MergeableTODS}, fix
$\varepsilon$ before the merge process begins. A summary method is allowed to
be one-to-many: for a dataset or stream multiset $D$, the expression
\[
S(D,\varepsilon)
\]
denotes any valid $\varepsilon$-summary that the method may output for $D$.
The size profile is
\[
k(n,\varepsilon)
  :=
  \max\{\, |S(D,\varepsilon)| : |D|=n,\ S(D,\varepsilon)\text{ valid}\,\}.
\]
The operation \(D_1\uplus D_2\) is multiset addition of represented inputs.
It is not addition, averaging, or merging of scalar query answers. Full
mergeability at the fixed $\varepsilon$ means that there is an algorithm $A$
such that, for every pair of valid child summaries,
\[
A\big(S(D_1,\varepsilon),S(D_2,\varepsilon)\big)
  \in S(D_1\uplus D_2,\varepsilon),
\]
and the resulting representation has size at most
\[
k(|D_1|+|D_2|,\varepsilon).
\]
If this profile is independent of $n$, we write it as $k(\varepsilon)$, so
both child summaries and the merged parent summary have size at most
$k(\varepsilon)$. The streaming case is the special case
\(|D_2|=1\), where each new item is merged into the current summary. Full
mergeability is stronger because the merge tree and the child dataset sizes are
arbitrary.

For randomized methods, the same notation is interpreted as a validity event:
after any finite merge tree, the root state must be a valid
$S(D,\varepsilon)$ for the represented multiset $D$ with the stated success
probability. C-TreePO's randomized nesting theorem preserves exactly this
event-level probability: on the event that the root state is a valid
$S(D,\varepsilon)$, the root readout is correct. The corresponding Lean
interface is \texttt{RandomizedRelationalMergeablePreferenceShape}; the
deterministic relational interface is \texttt{RelationalMergeablePreferenceShape}.

For the range-space examples, Agarwal's subset notation is also explicit. A
finite range space is \((D,\mathcal{R})\), and an
$\varepsilon$-approximation is a subset \(S\subseteq D\) satisfying
\[
\sup_{R\in\mathcal{R}}
\left|
  \frac{|S\cap R|}{|S|}
  -
  \frac{|D\cap R|}{|D|}
\right|
\le \varepsilon .
\]
In the C-TreePO nesting statement this subset \(S\) is the state/readout
object at the root, while \(D_1\uplus D_2\) is the represented-data operation
tracked by the tree. Keeping these two symbols separate prevents the common
mistake of treating set/subset readout notation as a law for merging scalar
child answers.

\paragraph{Conversion checklist.}
To convert a C-TreePO construction into Agarwal's interface, one must supply
the following objects, in this order:
\begin{enumerate}
\item a stream representation \(D\mapsto \mathrm{Stream}\,\alpha\) for the
represented data;
\item a fixed error parameter \(\varepsilon\);
\item a state type \(S\), a leaf builder
\(\mathrm{build}:\mathrm{Stream}\,\alpha\to S\), and a state merge
\(\mathrm{merge}:S\times S\to S\);
\item a validity relation \(V_\varepsilon(D,s)\), read as
``\(s\in S(D,\varepsilon)\)'';
\item a state-size function \(\mathrm{size}:S\to\mathbb{N}\) and a size profile
\(k(n,\varepsilon)\);
\item a root readout \(\mathrm{readout}:S\to Y\) and target oracle
\(F:\mathrm{Stream}\,\alpha\to Y\);
\item proofs that build produces valid states, merge preserves validity for
\(D_1\uplus D_2\), valid states respect the size profile, and valid states read
out to \(F(D)\).
\end{enumerate}
These are exactly the fields of the Lean adapter
\texttt{CTreePOToAgarwalTransform}. Its methods then produce the Agarwal
state-level summary interface, the C-TreePO relational nesting shape, the
leaf and merge valid-sized states, the arbitrary-tree valid-sized root state,
the root readout theorem, and the \(k(|D|,\varepsilon)\) size bound.

\subsection{Transformer Inclusion}

The transformer inclusion is architecture-level. Following Kovachki
Proposition~6, a pre-normalized single-head attention block agrees with
a discretized kernel layer on the block's active coordinates; stacking
attention, residual/local maps, and pointwise feed-forward layers
produces an equation-(6) neural operator. The membership statements we
use are therefore: (i) single-head attention equals a discretized kernel
layer, (ii) a transformer block is a neural-operator layer, (iii) a
finite transformer encoder stack is an equation-(6) neural operator,
and (iv) the encoder stack belongs to the equation-(6) class. These
membership claims are what let transformer-style summarizers use the
same external approximation-and-transfer route as the FNO examples:
equation-(6) approximation first, local-law budgets second, sampled
audit and method transport last. They operate at the architecture level
only; they do not certify tokenizers, runtime serving details, or a
specific model checkpoint.

\subsection{Mergeable-Sketch Inclusion}

The classical lane is broader than HLL. A sketch supplies an encoder, a state
merge, and a decoder/readout. HLL is the empirical parity anchor because it
gives a clean byte-identity test, but the theorem surface is the general
mergeable-sketch class, including the Lean-backed Count-Min,
Misra--Gries/GK-style, interval/range, bigram, and Markov-count interfaces.

The key class is $\mathrm{MS}(\fstar)$. A realized operator $g$ belongs to it
when it is induced by a sketch whose leaf state preserves the oracle value,
whose merge state preserves the oracle value across siblings, and whose
decoder/readout is compatible with the oracle. In that situation the local
laws are not an extra mystery: C1 is leaf validity, C3 is merge validity, and
C2 is stability when a stored summary is reused. The Lean theorem
\texttt{local\_laws\_of\_sketch} packages this direction, and
\texttt{mergeableSketchSummaryClass\_subset\_exactLocalLawSubspace} states the
corresponding neural-operator subspace membership.

This is a membership condition, not an automatic consequence of being a neural
operator. The covered class is exactly
$\mathcal{N}\cap \mathrm{MS}(\fstar)$, or an explicitly certified approximate
neighborhood of that class. In the calibrated-readout route below, the same
statement is read with $\widehat f$ in place of $\fstar$, and the separate
calibration theorem supplies the true-oracle slack.

\subsection{Universal Approximation Versus Certification}

The informal slogan ``neural operators can capture mergeable sketches'' should
be read as a two-step statement. First, neural-operator approximation gives
representational capacity. Fix a mergeable-summary interface
\[
\mathrm{build}:D\to S,\qquad
\mathrm{merge}:S\times S\to S,\qquad
\mathrm{readout}:S\to Y,
\]
and fix an encoding of the bounded state space $S$ into the operator's
input/output domain. On any compact realized-call set generated by a fixed
tree, a neural-operator class with the relevant approximation property can
approximate the leaf builder, the state merge, and the readout-facing state
map. This is the representation claim: the learned family can carry the same
kind of global state that the classical sketch carries on those calls.

Second, C-TreePO supplies the certification layer. Approximation alone does
not imply mergeability, and it does not certify a trained checkpoint. The
realized operator must satisfy the local laws on the state interface: C1 checks
that leaf calls build a state sufficient for the readout, C3 checks that
merging child states matches the raw-union state, and C2 checks that
already-produced states are stable under re-use. When these law budgets are
zero, the realized neural operator is in the exact state-level
mergeable-sketch subspace. When they are nonzero, the C-TreePO preservation
theorem gives a root distortion bound with the same local-law budget. With a
calibrated readout $\widehat f$ satisfying
\(\sup_x d(\widehat f(x),\fstar(x))\le\varepsilon_f\), the Lean bridge
adds only the two-sided calibration slack \(2\varepsilon_f\).

Thus the precise claim is:
\[
\text{NO capacity} + \text{state encoding} + \text{C1/C3/C2 certification}
\Longrightarrow
\text{learned mergeable-sketch behavior}.
\]
Optimizer convergence is a separate question. The local-law-weighted training
objective is best understood as a projection target: it tries to move the
realized neural operator into, or near, the mergeable-sketch subspace that the
theorem already certifies.

\subsection{Certified Overlap}

For a chosen neural-operator class $\mathcal{N}$, the certified overlap is
\[
\mathrm{MNO}_{\mathcal{N},\fstar}
  \;:=\;
  \mathcal{N} \cap \mathrm{MS}(\fstar),
\]
and this overlap sits inside the exact local-law neural-operator subspace
$\mathrm{ExactLL}_{\mathcal{N}}(T,\fstar)$. Separately,
Proposition~\ref{prop:mergeable-reduction}'s state-level route gives the
reverse direction of the qualified equivalence: an exact C-TreePO operator can
be read as a mergeable state scheme on the range of \(g\), with
\(\oplus_g(s,t)=g(s\concat t)\) and readout \(f\). This reverse direction is
state-level. It still does not assert that scalar child answers merge.

\subsection{Oracle \texorpdfstring{$+$}{+} Projection Decomposition}\label{app:oracle-plus-projection}

The operational learned setting usually lands in the approximate, not
exact, subspace. The formal objective therefore separates two roles:
calibrated readout fit and projection toward local-law validity. We write
$\widehat f$ for the learned readout calibrated against the true oracle
$\fstar$, and $g$ for the learned operator that builds leaf states and
merges child states. The local-law losses are evaluated relative to
$\widehat f$: C1 trains leaf-state sufficiency, C3 trains merge
consistency, and C2 trains stability on already-produced states. The
theorem-backed route is then:
\[
  \text{calibrate }\widehat f\approx \fstar,\qquad
  \text{project }g\text{ toward the C1/C3/C2 subspace for }\widehat f,
\]
and the Lean theorem
\begin{center}
\small
\path|trueOracle_delta_R_ZR_le_of_calibrated_neuralOperatorBridge|
\end{center}
transfers the resulting root distortion to $\fstar$ with additive
$2\varepsilon_f$ calibration slack.

A training objective matching that decomposition is
\begin{equation}\label{eq:calibrated-oracle-projection-objective}
\begin{aligned}
J(\widehat f,g)
  ={}& L_{\mathrm{cal}}(\widehat f,\fstar)
      + L_{\mathrm{root}}(\widehat f,g) \\
     &+ \lambda_{\mathrm{leaf}} C1_{\widehat f}(g)
      + \lambda_{\mathrm{merge}} C3_{\widehat f}(g)
      + \lambda_{\mathrm{idemp}} C2_{\widehat f}(g),
\end{aligned}
\end{equation}
with optional state-imitation terms in simulations where exact sketch
states are observed. This objective is a certification target, not an
optimizer theorem: it does not assert that SGD finds the projection point
without a separate optimization argument.

The main-text equation~\eqref{eq:oracle-projection-objective} writes the
projection part with a single scalar $\lambda\in[0,1]$, weighting the
three audited budgets uniformly. A more general per-law-weighted form
with non-negative shares $\rho_1,\rho_2,\rho_3$ reads
\begin{equation}\label{eq:oracle-projection-objective-generalized}
J_\lambda^{\boldsymbol{\rho}}(g)
  =
  (1-\lambda)\,L_{\mathrm{oracle}}(g)
  +
  \lambda\big(
    \rho_1\varepsilon_{\mathrm{leaf}}
    +\rho_2\varepsilon_{\mathrm{merge}}
    +\rho_3\varepsilon_{\mathrm{idemp}}\big).
\end{equation}
Setting $\rho_1=\rho_2=\rho_3=1$ recovers the main-text
$J_\lambda$; setting $\lambda=0$ gives pure oracle risk; setting
$\lambda=1$ gives the pure local-law projection term (for any fixed
share vector). The overall weight vector applied to the oracle term and
the three per-law terms is
$(1-\lambda,\,\lambda\rho_1,\,\lambda\rho_2,\,\lambda\rho_3)$, and the
objective bound flows through
\[
\text{Section 9 realization}
\;\to\;
\text{approximate local laws}
\;\to\;
\Delta_R \text{ / preference-gap bounds}.
\]
A fully-general per-budget weight vector $\boldsymbol{\lambda}=(\lambda_{\mathrm{leaf}},\lambda_{\mathrm{merge}},\lambda_{\mathrm{idemp}})$
that dispenses with the scalar/share decomposition subsumes the above
(with $\lambda_i := \lambda\rho_i$); we do not adopt it in the main text
because the scalar $\lambda\in[0,1]$ makes the iff endpoint theorems
cleaner to state.

This is the function-space reading of the training knob. Small
$\lambda$ prioritizes approximating the task oracle; large $\lambda$
prioritizes moving the realized operator into the local-law-constrained
subspace. The audit supplies the empirical certificate for where the
trained operator actually lands.

\subsection{Simplicity: Local-Law Weights Are a Projection}\label{app:projection-interpretation}

It is worth spelling out how simple the construction is when viewed
through the function-space lens. Fix a neural-operator class $\mathcal{N}$
--- the image of whatever parametric family is being trained --- and fix
an oracle $\fstar$ and tree $T$. The \emph{learnable local-law subspace}
is then just the intersection
\[
\mathrm{ExactLL}_{\mathcal{N}}(T,\fstar)
  \;:=\;
  \mathcal{N}\cap\mathrm{ExactLL}(T,\fstar),
\]
i.e.\ the operators that the class can represent \emph{and} that the
local laws permit. This is the same symbol already used in
Section~\ref{sec:theorems}. There is nothing in the training objective
that goes beyond pushing the learned operator into this set.

To make this precise, we separate two readings that are often conflated
in prose. The first --- ``local-law weights are a projection'' --- is
the set-level statement that the zero set of the local-law penalty is
exactly the exact local-law subspace. The second --- ``the
approximation error is structured by the local laws'' --- is the
pointwise statement that the only operators $g$ at which
$\mathrm{projectionPenalty}(g) = 0$ are the operators satisfying
C1/C2/C3. The following theorem records that these are the same
statement.

\begin{theorem}[Projection/structured-error equivalence]\label{thm:projection-iff}
Fix a tree $T$, an oracle $\fstar$, and a projection penalty
$\mathrm{projectionPenalty} : \Strings\rightsquigarrow\Strings\to\mathbb{R}_{\ge 0}$.
The following are equivalent:
\begin{enumerate}[leftmargin=1.5em,nosep]
\item \emph{(Projection framing.)}
$\{g : \mathrm{projectionPenalty}(g) = 0\} \;=\; \mathrm{ExactLL}(T,\fstar).$
\item \emph{(Structural-error framing.)} For every summarizer $g$,
$\mathrm{projectionPenalty}(g) = 0 \iff g \in \mathrm{ExactLL}(T,\fstar).$
\end{enumerate}
Moreover, for any fixed neural-operator class $\mathcal{N}$, both statements are
equivalent to
\[
\mathcal{N} \cap \{g : \mathrm{projectionPenalty}(g) = 0\}
  \;=\;
  \mathrm{ExactLL}_{\mathcal{N}}(T,\fstar).
\]
\end{theorem}

\noindent\emph{Proof.} Equivalence of (1) and (2) is set-extensionality
applied to the displayed sets. The class-restricted version follows by
intersecting both sides of (1) with $\mathcal{N}$ and using
$\mathcal{N}\cap\mathrm{ExactLL}(T,\fstar)=\mathrm{ExactLL}_{\mathcal{N}}(T,\fstar)$.
\qed

\medskip
In words: the local-law weights we add to the loss are, up to
faithfulness of the penalty, \emph{exactly} a projection operator onto
the learnable local-law subspace. The $\lambda\to 1$ limit of
$J_\lambda$ behaves like an indicator penalty for that set. Large
$\lambda$ is hard projection; small $\lambda$ is soft projection traded
off against oracle fit. The method is not an arbitrary regularizer; it
is a weighted penalty for leaving the subspace the theorems actually
need.

The operational $J_\lambda$ that the training loop uses is the
finite-sample realization of this picture with the \emph{approximate}
local-law bundle in place of an exact indicator. The audit of
\S\ref{sec:estimation} is what closes the loop: it reports where in
$\mathrm{ApproxLL}(T,\fstar;\varepsilon)$ the trained $g$ has landed,
and the bound in \S\ref{sec:discussion} converts that audited
$\varepsilon$ into the IPW term of Equation~\eqref{eq:discussion-bound}.

\paragraph{Projection framing $\iff$ structural assumption on the approximation error.}
Theorem~\ref{thm:projection-iff} is the precise content of the
equivalence the main text refers to as the projection/structured-error
iff. When we introduce $\lambda$-weighted local-law terms in training,
we are making a statement about what the trained operator is
approximating. Calling it a ``projection'' makes the \emph{target} (the
learnable local-law subspace) visible; calling it ``an assumption about
the structure of the approximation error'' makes the \emph{content of
the assumption} visible. They are the same assumption, and the iff is
what lets us move between them without sleight of hand.
